
\documentclass{article}
\usepackage{colortbl}
\usepackage{makecell}
\usepackage{multirow}
\usepackage{supertabular}

\begin{document}

\newcounter{utterance}

\centering \large Interaction Transcript for game `bbh', experiment `geometric\_shapes', episode 104 with Qwen3.6{-}27B.
\vspace{24pt}

{ \footnotesize  \setcounter{utterance}{1}
\setlength{\tabcolsep}{0pt}
\begin{supertabular}{c@{$\;$}|p{.15\linewidth}@{}p{.15\linewidth}p{.15\linewidth}p{.15\linewidth}p{.15\linewidth}p{.15\linewidth}}
    \# & \multicolumn{2}{c}{Player} && \multicolumn{2}{c}{Game Master} \\
    \hline

    \theutterance \stepcounter{utterance}  
    & & & \multicolumn{4}{p{0.6\linewidth}}{
        \cellcolor[rgb]{0.9,0.9,0.9}{
            \makecell[{{p{\linewidth}}}]{
                \texttt{\tiny{[P1$\langle$GM]}}
                \texttt{Name geometric shapes from their SVG paths.} \\
\\ 
\texttt{Q: This SVG path element <path d="M 31.00,73.00 L 32.00,59.00 L 44.00,50.00 L 49.00,41.00 L 64.00,37.00 L 71.00,55.00 L 64.00,76.00 L 52.00,61.00 L 31.00,73.00"/> draws a} \\
\texttt{Options:} \\
\texttt{(A) circle} \\
\texttt{(B) heptagon} \\
\texttt{(C) hexagon} \\
\texttt{(D) kite} \\
\texttt{(E) line} \\
\texttt{(F) octagon} \\
\texttt{(G) pentagon} \\
\texttt{(H) rectangle} \\
\texttt{(I) sector} \\
\texttt{(J) triangle} \\
\texttt{A: (F)} \\
\\ 
\texttt{Q: This SVG path element <path d="M 14.19,26.04 L 51.43,39.21 L 58.44,36.69 L 56.63,30.17 L 48.53,26.66 L 14.19,26.04"/> draws a} \\
\texttt{Options:} \\
\texttt{(A) circle} \\
\texttt{(B) heptagon} \\
\texttt{(C) hexagon} \\
\texttt{(D) kite} \\
\texttt{(E) line} \\
\texttt{(F) octagon} \\
\texttt{(G) pentagon} \\
\texttt{(H) rectangle} \\
\texttt{(I) sector} \\
\texttt{(J) triangle} \\
\texttt{A: (G)} \\
\\ 
\texttt{Q: This SVG path element <path d="M 41.00,43.00 L 37.00,34.00 L 41.00,33.00 L 45.00,34.00 L 41.00,43.00"/> draws a} \\
\texttt{Options:} \\
\texttt{(A) circle} \\
\texttt{(B) heptagon} \\
\texttt{(C) hexagon} \\
\texttt{(D) kite} \\
\texttt{(E) line} \\
\texttt{(F) octagon} \\
\texttt{(G) pentagon} \\
\texttt{(H) rectangle} \\
\texttt{(I) sector} \\
\texttt{(J) triangle} \\
\texttt{A: (D)} \\
\\ 
\texttt{Q: This SVG path element <path d="M 23.66,12.78 A 9.08,9.08 335.96 1,0 7.07,20.19 A 9.08,9.08 335.96 1,0 23.66,12.78"/> draws a} \\
\texttt{Options:} \\
\texttt{(A) circle} \\
\texttt{(B) heptagon} \\
\texttt{(C) hexagon} \\
\texttt{(D) kite} \\
\texttt{(E) line} \\
\texttt{(F) octagon} \\
\texttt{(G) pentagon} \\
\texttt{(H) rectangle} \\
\texttt{(I) sector} \\
\texttt{(J) triangle} \\
\texttt{(K) ellipse} \\
\texttt{A:} \\
            }
        }
    }
     \\ \\

    \theutterance \stepcounter{utterance}  
    & \multicolumn{4}{p{0.6\linewidth}}{
        \cellcolor[rgb]{0.561,0.737,0.561}{
            \makecell[{{p{\linewidth}}}]{
                \texttt{\tiny{[P1$\rangle$GM]}}
                \texttt{To determine the shape drawn by the SVG path, let's analyze the commands in the `d` attribute:} \\
\\ 
\texttt{`M 23.66,12.78`} \\
\texttt{This moves the pen to the starting point \$(23.66, 12.78)\$.} \\
\\ 
\texttt{`A 9.08,9.08 335.96 1,0 7.07,20.19`} \\
\texttt{This is an elliptical arc command.} \\
\texttt{{-} `rx=9.08`, `ry=9.08`: The radii are equal, indicating a circular arc.} \\
\texttt{{-} `x{-}axis{-}rotation=335.96`: The rotation of the ellipse.} \\
\texttt{{-} `large{-}arc{-}flag=1`: This indicates that the arc should sweep more than 180 degrees.} \\
\texttt{{-} `sweep{-}flag=0`: This indicates a counter{-}clockwise sweep.} \\
\texttt{{-} `x=7.07, y=20.19`: The end point of the arc.} \\
\\ 
\texttt{`A 9.08,9.08 335.96 1,0 23.66,12.78`} \\
\texttt{This is another elliptical arc command returning to the start point.} \\
\texttt{{-} Same radii and rotation.} \\
\texttt{{-} `large{-}arc{-}flag=1`: Again, sweeping more than 1} \\
            }
        }
    }
    & & \\ \\

    \theutterance \stepcounter{utterance}  
    & & & \multicolumn{2}{p{0.3\linewidth}}{
        \cellcolor[rgb]{0.9,0.9,0.9}{
            \makecell[{{p{\linewidth}}}]{
                \texttt{\tiny{[GM$|$GM]}}
                \texttt{Target: (K)} \\
            }
        }
    }
    & & \\ \\

    \theutterance \stepcounter{utterance}  
    & & & \multicolumn{2}{p{0.3\linewidth}}{
        \cellcolor[rgb]{0.9,0.9,0.9}{
            \makecell[{{p{\linewidth}}}]{
                \texttt{\tiny{[GM$|$GM]}}
                \texttt{game\_result = LOSE} \\
            }
        }
    }
    & & \\ \\

\end{supertabular}
}

\end{document}
